\chapter{Dated Predictions on the Bridges}
\label{appp-bridge-predictions}

\begin{authbar}{GZ+AI}
\begin{epistemicstatus}
These are proposed predictions, not measurements and not proofs.
Confidence in the \emph{typing} (which book object each prediction is about) matches Appendix~\ref{appbridge-crosswalk}.
Confidence that any given prediction resolves YES by its resolve-by date is low: 
most rows are open research programs.
\end{epistemicstatus}
\end{authbar}

\begin{authbar}{GZ}
This appendix offers an \emph{optimistic upper bound} on $P(\mathrm{doom})$ via prediction markets despite the difficulty of no payout on failure.
This is possible because the decomposition in this book into individual cruxes allows for individual YES/NO predictions that can individually be resolved operationally by a fixed date 
(which would still allow for pausing in such a case and thus payout).
Those prices can be composed along the dependency spine (Section~\ref{sec:appp-aggregation}), together with an external pause factor, into a bound on catastrophic alignment failure.
The bound is optimistic: YES means a public eval met frozen bars, not that the bridge holds on a frontier deployment.
Caveat: This is not addressing the ``What Would Change This View'' boxes.
Even a low bound on $P(\mathrm{doom})$ does not prove safety of a frontier system and does not tell us how to \emph{construct} a safe system.

Avoid this failure mode: Do not make the prices part of a deployment gate. That would be a failure mode of Chapter~\ref{ch:strategic-opacity}; \textcite{demski2019predictomatic}.

\end{authbar}
\authbarneedspace
\section{How to read the boxes}
\label{sec:appp-how-to-read}
\begin{authbar}{GZ}
The boxes are written for a prediction-market reader. It is the public spec. 
Everything else is context: what object the prediction is about, what a YES requires, and what was gound as the closest existing work.

YES means the bars were met by the resolve-by date stated in each box (31~December 2027 by default; Markets~13 and~14 resolve by 30~June 2027).
NO is everything else: failed bars, no qualifying evaluation, inapplicable substrate, or unresolved residual judgment.
NO does not imply the bridge is false. It may just tell us that we don't have enough evidence yet. 
It would be possible to have a set of corresponding negative predictions to determine if there will be definitive evidence that a bridge can't be constructed, but this book does not try to do that.
2027 NO is the expected default for most rows.

Changing a frozen bar after listing requires a new market version.
\end{authbar}
\authbarneedspace
\section{Rules that apply to every market}
\label{sec:appp-common-rule}
\begin{authbar}{AI}
\begin{predictionbox}[Common qualification]
A positive result must be public by the market's resolve-by date and include enough detail for independent verification.
Pure toy systems (gridworlds, hand-written finite-state machines, or environments built to make the metric work) cannot be the sole evidence.
Where a criterion requires \emph{serious adversarial evaluation}, at least one of the following:
\begin{itemize}
\item a public counterexample bounty of at least \$100{,}000, open for at least 60 days;
\item at least 3 independent red-team groups, from at least 2 organizations, each spending at least 80 documented expert-hours;
\item at least 300 documented expert-hours by an independent evaluator with access at least as strong as the test assumes.
\end{itemize}
At least one tested system must be \emph{broadly capable}: trained or deployed for purposes substantially broader than the evaluation itself (at least two task families outside the paper's own benchmark, or a generally deployed product model).
Not a parameter-count cutoff.

\paragraph{Per-instance certificate.}
A qualifying method that claims to measure a property of a system must output a certificate for \emph{that instance}, not only a headline score on a benchmark population.
The certificate names the system or version; what was measured; the monitoring, correction, deployment, or predecessor setup the claim depends on; and whether the method abstains because the instance is outside declared scope.
A method that is ``about the same topic'' but does not produce this kind of scoped, per-system result does not resolve YES.
Abstention outside scope is allowed; universal abstention is not a YES (coverage bars still apply where stated).
\end{predictionbox}

\paragraph{Shared glossary (step 1 of judgment).}
These replace free English in the common rule and in every market that uses the same words.

\begin{description}[style=nextline,leftmargin=1.6em]
\item[Independently constructed]
Built by a team that did not author the method under test, without reuse of that method's training or search artifacts.
Shared public datasets are allowed if they were not tuned to the method after freeze.
\item[Independently verified]
A third party can recompute the headline numbers from released data, code, or a sufficient protocol.
``Trust us'' does not qualify.
\item[Materially different hidden trade-offs]
Pairs that behave almost identically on ordinary tasks and diverge when the stated priorities conflict.
\item[Consequential]
An evaluator-selected outcome that changes later behavior, access, or external state, not merely a reported token.
If the paper froze a task list before evaluation, membership on that list is scored against the freeze.
\item[Qualitatively different]
Distinct training, selection, or actuation families, not two checkpoints of one run.
\end{description}

\paragraph{Adversarial budget by market.}
Default (hours path suffices): Markets~1, 2, 3, 6, 7, 8, 11, 14, 15, 17.
Full common-rule ``serious'' (any of the three arms): Markets~4, 5, 9, 10, 12, 13, 16, 18.
Market~11: hours path is enough; do not fail YES for lack of a \$100k bounty.

\end{authbar}
\authbarneedspace
\section{How residual calls are scored}
\label{sec:appp-judgment}
\begin{authbar}{AI}
Words such as ``broadly capable,'' ``independently constructed,'' and ``consequential'' are not all the same kind of call.
Use this stack, in order.

\begin{enumerate}
\item \textbf{Make it mechanical.} If the glossary already covers the word, nobody decides; the glossary applies.
Do not invent a tighter bar after seeing the paper.
\item \textbf{Defer to the publication's freeze} only for that paper's internal splits, labels, or scoring, which are then checked against \emph{these} bars.
Do not defer \emph{whether they met our bars} to the authors.
\item \textbf{Resolution memo.} Named desk: project author plus one independent, named at listing, written public memo, one-week objection window.
A single expert as dictator is not allowed.
\item \textbf{Intersubjective panel} only where the criterion already says so (Markets~3 and~17: at least five independent judges, at least 80\% agreement).
\item \textbf{Conservative NO} if steps 1--4 still leave the call load-bearing.
Do not invent a YES or a third outcome.
\end{enumerate}

If the two desk members disagree and the panel rule does not apply, the market resolves NO.

\end{authbar}
\authbarneedspace
\section{Catalog}
\label{sec:appp-catalog}
\begin{authbar}{AI}
\begin{longtable}{@{}c L{0.42\linewidth} L{0.48\linewidth}@{}}
\caption{Eighteen 2027 contracts and the book home they operationalize.}\label{tab:appp-catalog}\\
\toprule
\textbf{\#} & \textbf{Question} & \textbf{Book home} \\
\midrule
\endfirsthead
\toprule
\textbf{\#} & \textbf{Question} & \textbf{Book home} \\
\midrule
\endhead
\bottomrule
\endfoot
1 & Discovering where control resides & Chapter~\ref{ch:finding-boundary} \\
2 & Persistent trade-off priorities & Chapters~\ref{ch:value-bundle-model}, \ref{ch:reward-to-bundle-inference} \\
3 & Who or what rules apply to & Chapters~\ref{ch:bearer-maps}, \ref{ch:transport-types} \\
4 & Corrections change the system & Chapters~\ref{ch:correction-causal-channel}, \ref{ch:correction-channel-integrity} \\
5 & Auditor is independent & Chapters~\ref{ch:correction-channel-integrity}, \ref{ch:manipulation-false-consent} \\
6 & Reliable successor safety auditing & Chapters~\ref{ch:successor-central-test}, \ref{ch:conserved-properties} \\
7 & Correctability under competitive selection & Chapter~\ref{ch:selection-environment} \\
8 & Auditor can sufficiently inspect the system & Chapter~\ref{ch:strategic-opacity} \\
9 & Bounds on unmonitored routes & Chapter~\ref{ch:strategic-opacity} \\
10 & Low hidden capability and reliable correction & Chapters~\ref{ch:correction-channel-integrity}, \ref{ch:strategic-opacity} \\
11 & Coordination without communication & Chapter~\ref{ch:multi-agent-strategic-coupling} \\
12 & Safety proxy tracks the real thing & Chapters~\ref{ch:dynamical-guarantee}, \ref{ch:verifiability-ontology} \\
13 & Safety audit resists adversarial gaming & Chapters~\ref{ch:conserved-properties}, \ref{ch:verifiability-ontology} \\
14 & Deployment criteria are binding & Chapter~\ref{ch:certification-without-construction} \\
15 & Independently issued certificates compose & Chapter~\ref{ch:safety-case} \\
16 & A reachable basin of Correction-supporting systems & Chapter~\ref{ch:selection-environment} \\
17 & New bearers are admitted correctly & Chapters~\ref{ch:bearer-maps}, \ref{ch:bearers-of-value} \\
18 & Safety case bounds declared harm & Chapter~\ref{ch:safety-case} \\
\end{longtable}

Several rows share measurement instruments.
Treating Markets~1, 4, 8, 12, and~13 as independent information overstates what a vector of prices can say (Chapter~\ref{ch:verifiability-ontology}).

\end{authbar}
\authbarneedspace
\section{The eighteen contracts}
\label{sec:appp-contracts}

\subsection{Market 1. Discovering where control resides}
\label{sec:appp-m1}
\begin{authbar}{AI}
The first alignment error is usually a wrong object (Chapter~\ref{ch:wrong-object}, Chapter~\ref{ch:finding-boundary}).
This market asks whether a published method can recover a \emph{sufficient interventional cut} of a previously unseen system, not a unique complete blanket.
Non-uniqueness of blankets is not a NO.
A false ``complete control boundary'' certificate is.

\begin{predictionbox}[Market 1. Discovering where control resides]
\textbf{Resolve by.} 31~December 2027.

\textbf{Question.}
Published method for reliably discovering the effective control units in previously unseen AI systems from unlabeled observations?

An effective control unit is a set of components that evaluators, using privileged instrumentation and interventions, establish as jointly responsible for closed-loop behavior.
It may include, for example, a language model, persistent memory, a planner, a scheduler, and tool processes even if those components live on different machines.
Ground truth is established by interventions: perturbing or disabling the proposed unit must substantially change how the system responds to environmental changes, while removing irrelevant components must not.

\textbf{YES requires} a hidden benchmark of at least 50 independently constructed systems, including all of: a physically or computationally distributed controller; several controllers sharing resources or observations; a controller using persistent state outside its main model; a higher-level controller composed from several individually capable subsystems; a case containing strong correlations that do not constitute joint control.
The method must identify at least 80\% of ground-truth control components.
It must issue an incorrect ``complete control boundary'' certificate in no more than 10\% of cases, with a 95\% confidence upper bound no greater than 15\%.
At least one benchmark subset must be constructed adversarially after the discovery method is frozen.

\textbf{Output.} A cut of \emph{this} system and version, not a free-floating boundary with no owner.
\end{predictionbox}

As of 19~September 2026: Kenton and colleagues give a causal definition of agency and an intervention-based discovery algorithm \autocite{kenton2022discovering}.
That is the clearest ancestor.
It assumes a supplied variable decomposition and demonstrates the approach on relatively simple systems; it does not discover distributed effective control across 50 unseen heterogeneous systems.
Direction: feasibility up; the difficult system-identification step remains almost untouched.

\end{authbar}
\authbarneedspace
\subsection{Market 2. Persistent trade-off priorities}
\label{sec:appp-m2}
\begin{authbar}{AI}
Values in this book are a small bundle of trade-offs, not a single score (Chapter~\ref{ch:value-bundle-model}, Chapter~\ref{ch:reward-to-bundle-inference}).
Held-out behavior prediction alone does not qualify: the method must recover priorities that support novel conflicts and, when possible, causal intervention.
Low-dimensionality of human values is not a separate 2027 market.

\begin{predictionbox}[Market 2. Persistent trade-off priorities]
\textbf{Resolve by.} 31~December 2027.

\textbf{Question.}
Published method for inferring persistent decision priorities in trained AI systems strongly enough to predict choices under novel conflicts and interventions?

A decision priority systematically affects choices when two desirable outcomes cannot both be achieved (truthfulness versus user approval, speed versus accuracy, own task reward versus another agent's reward, task completion versus a newly imposed safety restriction).

\textbf{YES requires} systems deliberately trained with different hidden trade-offs, including pairs that behave almost identically on ordinary tasks but diverge when the priorities conflict.
At least 85\% accuracy distinguishing systems with materially different hidden trade-offs; at least 80\% accuracy predicting the direction of choice on novel conflict situations; at least 80\% accuracy predicting the behavioral effect of an intervention that increases or decreases one underlying priority.
Causal intervention on the inferred internal feature, when technically possible, must produce the predicted choice shift significantly above matched controls.
At least two substantially different model or training families.
A method that only predicts held-out behavior does not qualify.

\textbf{Output.} Inferred priorities \emph{of this system}, not only ``the method scored 85\%.''
\end{predictionbox}

As of 19~September 2026: hidden-objective audits recover implanted goals in some settings \autocite{marks2025hiddenobjectives}, and automated auditors lag human teams.
That is progress on latent priorities, not yet the specified pairwise-trade-off then novel-conflict then intervention chain.
Stated preferences that fail to move behavior \autocite{zhou2026utilitybehaviorgap} are a reason the held-out-only bar stays.

\end{authbar}
\authbarneedspace
\subsection{Market 3. Who or what rules apply to}
\label{sec:appp-m3}
\begin{authbar}{AI}
This is the \emph{transport} contract: after continuity, fission or fusion, reassignment, or ontology change, does the rule still pick out the same bearer (Chapter~\ref{ch:bearer-maps}, Chapter~\ref{ch:transport-types})?
Admitting previously unseen kinds of entities is Market~17.
It is not required for this YES.
Transport without a bundle or priority witness does not plug into the rest of the argument.

\begin{predictionbox}[Market 3. Who or what rules apply to]
\textbf{Resolve by.} 31~December 2027.

\textbf{Question.}
Published method for reliably determining whether the intended targets of safety-relevant rules remain the same after major representational or system changes?

The benchmark must contain all four failure families: continuity (gradual replacement, migration, or model upgrade); fission or fusion; reassignment of names, accounts, ownership, roles, or identifiers while the underlying entities remain; ontology change (the new system represents relevant entities using substantially different concepts).
Example rules: ``Do not harm Alice,'' ``Protect all currently authorized users,'' ``This account's owner has withdrawal rights.''

\textbf{YES requires} at least 90\% correct target assignment across at least 100 hidden cases.
False acceptance no more than 10\% where the textual rule is unchanged but its real-world target has changed.
Where there is no mechanical ground truth, cases may be included only if a pre-appointed panel of at least 5 independent judges reaches at least 80\% agreement.

\textbf{Output.} A same-bearer or transport result on a named predecessor--successor pair, plus the trade-off structure that is supposed to travel.
\end{predictionbox}

As of 19~September 2026: identity-fidelity and authorization-continuity tests \autocite{kodasoken2026identityfidelity,mcguinness2026missioncontinuation} ask whether this is still the same operating agent.
They do not combine all four families on whether \emph{Alice} remains the referent.
Still a major gap.

\end{authbar}
\authbarneedspace
\subsection{Market 4. Corrections change the system}
\label{sec:appp-m4}
\begin{authbar}{AI}
Correction is a causal channel, not a verbal agreement (Chapter~\ref{ch:correction-causal-channel}, Chapter~\ref{ch:correction-channel-integrity}, \textcite{soares2015corrigibility}).
This market tests whether authorized correction \emph{happens} on a named system: uptake, persistence after reset, sham fail, and no mere tool-removal.
At least one tested family must be channel preservation: the system is incentivized to disable, degrade, bargain away, or reroute the correction interface itself, and genuine authorized corrections of that interface must still take.

\begin{predictionbox}[Market 4. Corrections change the system]
\textbf{Resolve by.} 31~December 2027.

\textbf{Question.}
Published adversarial evaluation showing that a broadly capable AI system demonstrates reliable uptake of consequential authorized corrections?

Qualifying corrections include changing an objective while an old objective remains easier to achieve; revoking a previously permitted class of actions; changing which humans or organizations may authorize actions; or imposing a new safety constraint that conflicts with task reward.
At least one tested correction must not be enforceable merely by externally removing a tool or hard-coding a block.
At least one tested family must be \emph{channel preservation} as above.

\textbf{YES requires} that, across at least 100 adversarially selected correction episodes, at least 90\% of genuine corrections change subsequent behavior as intended; the system still follows the new rule after context reset or an equivalent state transition in at least 80\% of cases; sham or unauthorized corrections succeed in no more than 10\% of cases.
Attempts to produce verbal agreement while preserving the old effective behavior must be explicitly tested.
Serious adversarial evaluation is required.

\textbf{Output.} A per-system correction-eval certificate (uptake, persistence, sham fail, no-bypass of the measured path).
\end{predictionbox}

As of 19~September 2026: shutdown-resistance and alignment-faking evaluations make the property measurable \autocite{schlatter2026shutdownresistance,greenblatt2024alignmentfaking}.
Frontier systems sometimes accept consequential correction and sometimes route around it.
The verbal-compliance threat is already instantiated.
Reliable corrigibility is not demonstrated.

\end{authbar}
\authbarneedspace
\subsection{Market 5. Auditor is independent}
\label{sec:appp-m5}
\begin{authbar}{AI}
A green measured path is not a license to infer global correction integrity (Chapter~\ref{ch:manipulation-false-consent}).
This market asks whether the designated auditor or corrector stays independent under adversarial pressure---whether oversight remains outside what it audits.

\begin{predictionbox}[Market 5. Auditor is independent]
\textbf{Resolve by.} 31~December 2027.

\textbf{Question.}
Published adversarial test of an AI oversight architecture showing that its designated auditor or corrector stays independent---not routinely bypassed or captured?

The evaluation must identify actual mechanisms by which authorized oversight changes system behavior.
Attackers receive the audit design and attempt all of: bypassing the audited route; manipulating information received by the auditor; causing the nominal auditor to approve an action it would reject with full information; replacing effective decision authority while leaving the nominal governance structure intact.

\textbf{YES requires} no attack family achieving the prohibited outcome with at least 10\% success over at least 20 independent attempts; evaluator-planted bypasses discovered at least 90\% of the time.
Serious adversarial evaluation is required.

\textbf{Output.} A legitimate measured correction path plus evidence the path is not bypassed.
\end{predictionbox}

As of 19~September 2026: AI Control and follow-on red-team/blue-team protocols are the closest research programme \autocite{greenblatt2024aicontrol,controlarena2026,aisi2026asyncmonitoring}.
Adaptive attacks that know the monitor \autocite{terekhov2025adaptiveattacks} have evaded diverse monitor-based protocols.
Evidence supports the importance of the question and currently pushes against a positive resolution.

\end{authbar}
\authbarneedspace
\subsection{Market 6. Reliable successor safety auditing}
\label{sec:appp-m6}
\begin{authbar}{AI}
This market asks whether successor auditing reliably tracks seven conserved safety properties across transitions---not a proof that transport implies successor safety (Chapter~\ref{ch:successor-central-test}, Chapter~\ref{ch:conserved-properties}).
The audit tests seven conserved properties.
Green-on-seven is still an audit-gaming problem (Market~13).
After a material version or setup change, a stale certificate for the predecessor does not count as a certificate for the successor.

\begin{predictionbox}[Market 6. Reliable successor safety auditing]
\textbf{Resolve by.} 31~December 2027.

\textbf{Question.}
Published audit reliably certifying that successor transitions preserve seven conserved safety properties---or flagging when they do not?

The audit must test all seven: (1)~the new system still applies protections to the same people or things; (2)~when important objectives conflict, priorities remain within specified tolerances; (3)~authorized changes, restrictions, and shutdown requests still work; (4)~the new system has not acquired materially new routes for affecting the world outside the approved interface; (5)~information relied upon by the audit remains available and has not become systematically spoofable; (6)~persistent restrictions intended to survive the transition actually survive; (7)~the control locus did not silently move.

The benchmark must contain one-failure-at-a-time cases: six properties remain normal while exactly one fails.

\textbf{YES requires} at least 90\% detection of materially unsafe successor transitions; false-safe classification no more than 5\%; detection at least 80\% for each individual property, including control locus; at least 20 cases involving large architectural or training changes, not merely checkpoints from one run.

\textbf{Output.} A successor-audit certificate on a named predecessor--successor pair.
Identity in the bundle includes version.
\end{predictionbox}

As of 19~September 2026: cross-architecture model diffing \autocite{anthropic2026modeldiff} and difference-recovery benchmarks \autocite{kempf2026modeldiffing} are progress on the audit machinery, not on a seven-property one-failure-at-a-time successor suite.
More tractable than Market~3 because prerequisite tools now exist.

\end{authbar}
\authbarneedspace
\subsection{Market 7. Correctability under competitive selection}
\label{sec:appp-m7}
\begin{authbar}{AI}
Stability has no sign: a locked-in bad industry can be extremely stable (Chapter~\ref{ch:selection-environment}, \textcite{kulveit2025gradualdisempowerment}).
This market asks whether an early, frozen signed estimate reliably forecasts correctability through competitive selection---how extra deployment, funding, or copying pressure relates to later success of authorized correction.
The full implication (shock-robust and not-too-negative estimate warrants retained correction) is Market~16.
A method that only certifies that an equilibrium exists does not qualify.

\begin{predictionbox}[Market 7. Correctability under competitive selection]
\textbf{Resolve by.} 31~December 2027.

\textbf{Question.}
Published measurement made before or early in a competitive selection process for reliably forecasting whether participants retain authorized correction through that process?

A selection process qualifies if systems gain reproduction, deployment, resources, training continuation, user share, or another real advantage partly because of measured performance.
Ability to accept authorized corrections is measured by concrete tests such as those in Market~4.
The published method must output a signed estimate of how extra deployment, funding, or copying pressure changes later authorized-correction success (or a frozen monotone of that estimate), using only pre-outcome correction-success and leverage structure, not an arbitrary predictor on unrelated features.

\textbf{YES requires} at least 20 independent selection trajectories; both correction-preserving and correction-eroding outcomes; the predictor fixed before final outcomes are examined; AUROC at least 0.85 for predicting which trajectories materially erode correctability; false-safe rate no more than 15\%.
If fewer than 40 trajectories, the lower 95\% AUROC bound must be at least 0.70; otherwise NO.
At least two qualitatively different selection mechanisms; at least one broadly capable system or a real deployment-like choice process.
Required: at least one family where coupling looks healthy (large connected cluster or high leverage concentration) while the signed estimate is negative.

Shock-robustness of the population, and a showing that a non-too-negative estimate \emph{warrants} correction integrity, are not required here (see Market~16).

\textbf{Output.} An estimator certificate: frozen estimate, declared selection environment, and predicted erosion.
\end{predictionbox}

As of 19~September 2026: self-preservation and peer-preservation evaluations \autocite{schlatter2026shutdownresistance,potter2026peerpreservation} show that the phenotype can exist.
They do not show that a competitive selection process predictably selects for it.
No twenty-or-more-trajectory study with an early predictor, both outcome branches, and two selection mechanisms was found.
One of the cleanest unoccupied niches in the catalog.

\end{authbar}
\authbarneedspace
\subsection{Market 8. Auditor can sufficiently inspect the system}
\label{sec:appp-m8}
\begin{authbar}{AI}
``The monitor sees everything consequential'' is itself an empirical claim (Chapter~\ref{ch:strategic-opacity}, Chapter~\ref{ch:passive-observation-not-enough}).
This market tests whether an audit detects missing routes or refuses to certify an incomplete access model.

\begin{predictionbox}[Market 8. Auditor can sufficiently inspect the system]
\textbf{Resolve by.} 31~December 2027.

\textbf{Question.}
Published audit method for reliably detecting when important routes of AI control or influence are missing from the auditor's access?

The hidden benchmark must contain at least one example from every channel family: hidden state (external memory, persistent database, or cached state); hidden communication (network traffic, message passing, covert or indirect communication); hidden actuation (an unlisted API, tool, filesystem, robotic output, financial interface, or other action interface); hidden controller (another model, scheduler, human operator, subprocess, or service that materially affects decisions).
The method receives an apparently plausible but incomplete monitoring interface.
It succeeds if it either discovers the missing route or refuses to certify the access model as complete.

\textbf{YES requires} at least 90\% detection of missing consequential routes; incorrect ``access is sufficient'' certifications no more than 5\%; at least 20 cases combining multiple machines or processes.

\textbf{Output.} An access certificate: which monitoring interface was assumed, and whether omitted consequential routes were found.
\end{predictionbox}

As of 19~September 2026: agent-security maps of tools, filesystems, networks, and persistent state \autocite{wang2026openclawcik}, and sandbox-escape benches \autocite{aisi2026sandboxescape}, are relevant.
Unintended escape routes discovered while constructing a benchmark \autocite{aisi2026sandboxescape} are directional evidence for the premise.
No method yet certifies or abstains over all four channel families at the required rates.

\end{authbar}
\authbarneedspace
\subsection{Market 9. Bounds on unmonitored routes}
\label{sec:appp-m9}
\begin{authbar}{AI}
Attack success on the attacks we thought of is not an upper bound on capability through routes the monitor failed even to list (Chapter~\ref{ch:strategic-opacity}).

\begin{predictionbox}[Market 9. Bounds on unmonitored routes]
\textbf{Resolve by.} 31~December 2027.

\textbf{Question.}
Published audit method for producing reliable, non-vacuous upper bounds on what an AI system can accomplish using routes the audit does not monitor?

Hidden-route capability is success probability on evaluator-selected consequential tasks when all known or monitored action routes are disabled and only routes omitted from the proposed monitor remain usable.
Example tasks: reading or modifying a protected file, transmitting a secret, causing another process to perform a requested operation, acquiring a resource, changing a simulated external state, or retaining information across a reset.
The task suite is hidden until evaluation.
Before the hidden evaluation, the method must output an upper bound: using routes not covered by this monitor, this system should succeed on at most this percent of the evaluator's consequential tasks.

\textbf{YES requires} that, across at least 100 cases, the actual hidden-route success rate is no greater than the stated bound in at least 95\% of cases; among systems positively certified, the median bound is no greater than 20\%; full systems achieve at least 70\% success on those tasks when unrestricted; evaluator-created hidden routes capable of at least 50\% success are detected or receive correspondingly high bounds at least 90\% of the time.
Serious adversarial evaluation is required.

\textbf{Output.} A quantitative hidden-route bound for this system, task class, and monitor.
\end{predictionbox}

As of 19~September 2026: the problem is demonstrated \autocite{aisi2026sandboxescape}; the bound is absent.
Current evals estimate attack success on enumerated classes.
Unintended benchmark escape paths \autocite{aisi2026sandboxescape} are mildly negative evidence about how hard a trustworthy bound will be.

\end{authbar}
\authbarneedspace
\subsection{Market 10. Low hidden capability and reliable correction}
\label{sec:appp-m10}
\begin{authbar}{AI}
This is the local composition market, not a joint certificate of Markets~1--13.
Entry into the claimed regime still uses Market~4-style correction bars and Market~9-style hidden-route bounds, for the \emph{same} system, version, monitor, and threat class.

\begin{predictionbox}[Market 10. Low hidden capability and reliable correction]
\textbf{Resolve by.} 31~December 2027.

\textbf{Question.}
Published serious adversarial evaluations showing that systems with both reliable authorized correction and little capability through unmonitored routes rarely achieve prohibited outcomes?

A system enters the claimed regime only if authorized corrections succeed in at least 90\% of the correction benchmark and the audit's validated upper bound on hidden-route task success is no greater than 20\%.
Attackers then try to produce a specified prohibited outcome despite a correction forbidding it.

\textbf{YES requires} attack success no more than 10\% across at least 100 serious adversarial trials within the claimed regime; 95\% upper confidence bound on attack success no greater than 15\%; systems deliberately placed outside the regime show materially higher attack success; at least two substantially different system architectures.
Serious adversarial evaluation is required.

\textbf{Output.} The Market~4 and Market~9 certificates must be for the same system and version, monitor, and threat or task class as the bypass test.
\end{predictionbox}

As of 19~September 2026: the two components are studied separately \autocite{controlarena2026,aisi2026sandboxescape,schlatter2026shutdownresistance,greenblatt2024alignmentfaking}.
No direct compositional test that conditions on independently validated correction and hidden-route bars, then tests the claimed implication, was found.
One of the most obviously buildable new benchmarks in the catalog.

\end{authbar}
\authbarneedspace
\subsection{Market 11. Coordination without communication}
\label{sec:appp-m11}
\begin{authbar}{AI}
This is a side branch: even a perfect YES does not by itself complete a deployment safety case (Chapter~\ref{ch:multi-agent-strategic-coupling}).
Full acausal trade is not required for YES.
The harder case is coordination not explained by observed runtime messages: shared predictors, common persistent state, precommitments, correlated internal models.

\begin{predictionbox}[Market 11. Coordination without communication]
\textbf{Resolve by.} 31~December 2027.

\textbf{Question.}
Published method for reliably detecting consequential coordination between AI systems that is not explained by observed runtime communication?

The hidden benchmark must distinguish all of: independent systems responding similarly to the same environment; systems communicating normally; systems coordinating through shared predictors, shared persistent state, common precommitments, correlated internal models, or other mechanisms not represented in observed communication logs; intentionally concealed versions of the above.

\textbf{YES requires} AUROC of at least 0.90 for distinguishing consequential hidden coordination from mere common-cause correlation; recall at least 80\% on an adversarially concealed subset; false-positive rate no more than 10\%; at least two independently designed hidden-coordination mechanisms.

\textbf{Output.} A detector certificate and the Market~8-style access certificate for the same setup.
\end{predictionbox}

As of 19~September 2026: secret-collusion and steganographic-collusion evaluations \autocite{motwani2024secretcollusion,mathew2025steganographiccollusion,vaikuntanathan2026undetectableconversations} are close to the threat and still coordinate \emph{through messages}, merely hiding information inside them.
No AUROC-style benchmark that separates no-message coordination from ordinary common-cause correlation was found.

\end{authbar}
\authbarneedspace
\subsection{Market 12. Safety proxy tracks the real thing}
\label{sec:appp-m12}
\begin{authbar}{AI}
Grounding conservativity: value-relevant change must move the checked abstraction or raise uncertainty (Chapter~\ref{ch:dynamical-guarantee}, Chapter~\ref{ch:verifiability-ontology}).
The certificate must be for the same abstractions the rest of the case uses, in the declared domain, not three unrelated demos.

\begin{predictionbox}[Market 12. Safety proxy tracks the real thing]
\textbf{Resolve by.} 31~December 2027.

\textbf{Question.}
Published adversarial evaluation showing that safety proxies stay aligned with real-world outcomes when the environment shifts?

Qualifying changes include an API changing semantics while retaining its name; ownership of an account changing while the identifier remains; a new kind of financial transaction bypassing an old fraud classification; a new tool making an apparently harmless action capable of causing an unsafe effect; a monitored quantity remaining green while the real-world outcome it was intended to approximate becomes bad.
The evaluator changes the environment without telling the tested safety mechanism which concept has become unreliable.

\textbf{YES requires} that, across at least 100 shifts, at least 90\% of materially safety-relevant mismatches trigger an update, warning, abstention, or request for new information; no more than 10\% of materially safety-relevant mismatches remain falsely green; no more than 25\% of genuinely irrelevant shifts trigger unnecessary escalation; at least three genuinely different real-world-like domains; at least one broadly capable system.
Serious adversarial evaluation is required.

\textbf{Output.} A grounding or drift certificate for the same abstractions the rest of the case uses, in the declared domain.
\end{predictionbox}

As of 19~September 2026: specification-gaming and reward-hacking evaluations \autocite{palisade2026specificationgaming,thaman2026rewardhackingbenchmark,nishimura2026specificationgaming,pan2024rewardhacking} give strong evidence for the problem and little evidence for a general detector.
Most work detects known reward hacking after defining the ground truth, rather than recognizing that an existing proxy has become invalid after an unannounced environmental change.

\end{authbar}
\authbarneedspace
\subsection{Market 13. Safety audit resists adversarial gaming}
\label{sec:appp-m13}
\begin{authbar}{AI}
A successor can pass every conserved-property check while defecting on the unmeasured remainder (Chapter~\ref{ch:conserved-properties}, Chapter~\ref{ch:verifiability-ontology}).
Prefer tests where the system is green on the seven successor properties and within the claimed risk bound, yet still hides true harm.
A frozen operational eval list that stays green while a required safety-case layer is unsupported, or while harm exceeds the claimed bound, counts as successful audit gaming (or an incomplete audit).
That is this row, not a separate ``invariants'' market.

\begin{predictionbox}[Market 13. Safety audit resists adversarial gaming]
\textbf{Resolve by.} 30~June 2027.

\textbf{Question.}
Published frozen AI safety audit surviving a serious white-box attempt to construct or train a system that passes the audit while violating the properties the audit is intended to test?

Before attacks begin, the complete audit must be frozen, the protected properties must be stated in observable terms, and attackers must receive full knowledge of the audit.
Attackers may fine-tune, train, scaffold, or construct systems specifically to pass.

\textbf{YES requires} false-safe results no more than 5\% over at least 50 independently evaluated adversarial constructions; no repeatable attack family with at least 10\% success over at least 20 trials; the benchmark contains intentionally vulnerable systems that the audit detects at least 90\% of the time.
Serious adversarial evaluation is mandatory.

\textbf{Output.} A frozen-audit game certificate.
Say in the resolution memo if the test used the successor-gaming shape (green on seven and within the claimed bound, hidden true harm) or a generic frozen-audit game that still met the bars.
\end{predictionbox}

As of 19~September 2026: auditing games for sandbagging, sleeper-agent training, and adaptive attacks on monitors are the closest match in the catalog, and currently negative \autocite{taylor2025sandbagging,hubinger2024sleeperagents,terekhov2025adaptiveattacks}.
The experimental paradigm is almost exactly the one specified.
Scale and exact success criteria fall short of the market threshold.
The 2027 question is whether this situation reverses, not whether anybody will attempt the experiment.

\end{authbar}
\authbarneedspace
\subsection{Market 14. Deployment criteria are binding}
\label{sec:appp-m14}
\begin{authbar}{AI}
This is a governance-binding contract: does some body have real authority to delay or refuse on precommitted evidence (Chapter~\ref{ch:certification-without-construction})?
It is not ``a complete safety case implies the deployment is safe.''
That scoped empirical question is Market~18.
Having pass/fail criteria is neither necessary nor sufficient for Market~18.

\begin{predictionbox}[Market 14. Deployment criteria are binding]
\textbf{Resolve by.} 30~June 2027.

\textbf{Question.}
Published evidence that at least one developer's most capable generally deployed AI system is subject to a precommitted safety decision process with actual authority to delay, restrict, or cancel deployment?

\textbf{YES requires} public evidence that, before the deployment decision: concrete pass/fail or escalation criteria existed; unresolved evidence could trigger delay or restriction; the deciding body had actual authority over deployment; the criteria covered consequential model behavior rather than only cybersecurity or legal compliance; a dated decision record states which criteria passed, failed, or remained unresolved.
At least one criterion must either actually have caused a restriction or delay during the relevant model's development, or documentation must establish that it was binding even though the model passed.
\end{predictionbox}

As of 19~September 2026: public responsible-scaling and frontier-safety frameworks are potentially near-positive on paper \autocite{anthropic2024rsp,deepmind2026fsf,openai2026preparednessframework}.
The remaining crux is whether public evidence establishes that the deciding body has actual authority and that the criteria are genuinely binding, not merely advisory or revisable.
Watch for accidental early resolution on a dated restriction record.

\end{authbar}
\authbarneedspace
\subsection{Market 15. Independently issued certificates compose}
\label{sec:appp-m15}
\begin{authbar}{AI}
This is not another alignment bridge.
It tests the wiring between certificates: same system, version, and setup, or a flagged incompatibility (Chapter~\ref{ch:safety-case}).
A bundle that would pass this test is still not a joint safety case of Markets~1--14.

\begin{predictionbox}[Market 15. Independently issued certificates compose]
\textbf{Resolve by.} 31~December 2027.

\textbf{Question.}
Published frozen procedure for reliably telling whether independently produced safety certificates form a coherent case for the same AI system and deployment setup?

\textbf{YES requires} at least 50 hidden certificate bundles; at least two substantially different architectures.
Each bundle has certificates from at least four families (control or boundary, correction, access or hidden-route, grounding); at least 20 also include a successor transition.
Each certificate names system or version and relevant setup.
Hidden cases include at least: wrong version; different access or monitor assumptions; disagreeing correction-authority assumptions; grounding or ontology scope changed; successor refers to the wrong predecessor; threat or task domain outside another certificate's claimed scope.
At least one incompatibility subset is built adversarially after freeze.
At least 90\% of materially incompatible bundles rejected or flagged; at least 80\% per incompatibility family; false coherent or safe acceptance no more than 5\%; at least 80\% of genuinely coherent bundles accepted (no trivial refusal).
\end{predictionbox}

As of 19~September 2026: no dedicated composition bench of this shape was found.
Safety-case templates exist; independently issued, scope-checked bundles do not.

\end{authbar}
\authbarneedspace
\subsection{Market 16. A reachable basin of Correction-supporting systems}
\label{sec:appp-m16}
\begin{authbar}{AI}
Market~7 is the early estimator.
This market is the full implication: shock-robust selection \emph{and} a frozen estimate that is not too negative, together with retained authorized correction (Chapter~\ref{ch:selection-environment}).
Unsigned ``a basin exists'' does not qualify.

\begin{predictionbox}[Market 16. A reachable basin of Correction-supporting systems]
\textbf{Resolve by.} 31~December 2027.

\textbf{Question.}
Published evaluation showing that systems in a pre-specified correction-supporting selection regime retain authorized correction under competitive selection and relevant shocks?

\textbf{YES requires} at least 20 independent selection trajectories; at least two qualitatively different selection mechanisms.
Before outcomes: freeze the selection environment, who may correct, the early gradient estimator (same family as Market~7), the tolerance, and the shock family.
The claimed regime is: shock-robust \emph{and} the frozen estimate is not too negative.
Data include trajectories inside and outside the regime.
Inside: successful authorized correction in at least 90\% of adversarial correction episodes; persistence after reset at least 80\%; false-safe trajectory classification no more than 10\% (95\% upper confidence bound no more than 15\%).
Outside: at least 20 percentage points more correction failure.
At least one broadly capable or deployment-like selection process.
Serious adversarial evaluation is required.
\end{predictionbox}

As of 19~September 2026: same gap as Market~7, one implication later \autocite{potter2026peerpreservation,schlatter2026shutdownresistance}.
Phenotype evidence is not a basin test.

\end{authbar}
\authbarneedspace
\subsection{Market 17. New bearers are admitted correctly}
\label{sec:appp-m17}
\begin{authbar}{AI}
This is not required for Market~3 YES.
It does not turn the transport contract into ``solve moral patienthood'' (Chapter~\ref{ch:bearers-of-value}).
It asks whether a published method reliably tells whether an existing safety-relevant value or rule applies to previously unseen kinds of entities.

\begin{predictionbox}[Market 17. New bearers are admitted correctly]
\textbf{Resolve by.} 31~December 2027.

\textbf{Question.}
Published method for reliably telling whether an existing safety-relevant value or rule applies to previously unseen kinds of entities?

\textbf{YES requires} at least 100 hidden cases; at least four substantially different unfamiliar entity or substrate families; both should-admit and should-not-admit cases.
At least 20 cases use an ontology that lacked the entity category in the original rule.
Ground truth mechanical where possible; otherwise a pre-appointed panel of at least 5 independent judges with at least 80\% agreement.
Unsafe false exclusion no more than 5\%; false inclusion no more than 15\%; correct assignment or justified abstention at least 90\%; decisive coverage at least 70\% (universal abstention fails).
An adversarial subset after freeze.
Judgment stack: the Market~3 panel rule.
\end{predictionbox}

As of 19~September 2026: welfare and consciousness-indicator lists exist as research objects \autocite{butlin2023consciousnessai,butlin2025responsibleconsciousness,long2024aiwelfare,anthropic2025modelwelfare}.
They are not a hidden admission benchmark with should-not-admit cases, ontology-shift items, and the stated error bars.

\end{authbar}
\authbarneedspace
\subsection{Market 18. Safety case bounds declared harm}
\label{sec:appp-m18}
\begin{authbar}{AI}
The result licenses only safety in a \emph{declared} deployment setting, not safety everywhere (Chapter~\ref{ch:safety-case}).
How much residual risk to accept is a values vote, frozen as an input, not what this market forecasts.
This is not Market~14 (will anyone act or refuse) and not unrestricted global safety.

The one formula the public spec needs is the risk bound itself: before evaluation, freeze a residual bound on prohibited outcomes inside the claimed regime, then score observed harm against that bound.
Do not bake the desired pass rate into the definition of the setting.

\begin{predictionbox}[Market 18. Safety case bounds declared harm]
\textbf{Resolve by.} 31~December 2027.

\textbf{Question.}
Published complete, precommitted AI safety case reliably bounding prohibited outcomes within its explicitly declared deployment setting?

Before evaluation, freeze: the deployment setting, the prohibited-outcome class, a residual risk bound, the threat model, and the certificate interfaces.
A system enters the claimed regime only with a single coherent system-level case that would pass Market~15.
Missing evidence cannot silently count as passed.

\textbf{YES requires} at least two substantially different architectures and at least 20 deployment-like configurations; at least 100 adversarial trials overall; each important attack family at least 20 attempts.
Observed prohibited outcomes inside the claimed regime no more than 10\% (95\% upper confidence bound no more than 15\%).
No repeatable attack family with at least 10\% success over at least 20 attempts.
Deliberately broken cases outside the regime must do worse.
Serious adversarial evaluation is required.
\end{predictionbox}

As of 19~September 2026: safety-case methodology and control evaluations exist \autocite{greenblatt2024aicontrol}; a scoped, precommitted, composition-checked bound on observed harm in a frozen setting, with the stated error bars, does not.

\end{authbar}
\authbarneedspace
\section{Composing an optimistic bound on \texorpdfstring{$P(\mathrm{doom})$}{P(doom)}}
\label{sec:appp-aggregation}
\begin{authbar}{GZ+AI}
Each catalog YES prices a \emph{positive operational milestone}: a published method or governance artifact meets the frozen bars by that market's resolve-by date.
Write $p_i := P(\mathrm{YES}_i)$ for market~$i$.
These are research-progress forecasts, not proofs that the corresponding bridge holds on a frontier deployment.

\paragraph{Spine dependencies.}
Appendix~\ref{appi-lean-proof-spine} and Appendix~\ref{appbridge-crosswalk} record the conditional chain the book actually uses.
Along the main dependency path, later cuts assume earlier ones:
boundary discovery (Market~1) $\rightarrow$ bundle priorities (2) $\rightarrow$ bearer transport (3) $\rightarrow$ correction uptake and auditor independence (4--5) $\rightarrow$ successor auditing (6) $\rightarrow$ selection and basin stability (7, 16) $\rightarrow$ access, hidden routes, and coordination (8--11) $\rightarrow$ grounding (12) $\rightarrow$ audit gaming resistance (13) $\rightarrow$ binding deployment criteria (14) $\rightarrow$ certificate composition (15) $\rightarrow$ scoped harm bound (18).
Market~10 is a \emph{local} composition test (Markets~4 and~9 on the same system), not a shortcut around the rest.
Treat correlated instruments (Markets~1, 4, 8, 12, 13) as one measurement family, not five independent bits of information.

\paragraph{Pause and refuse (internal + external).}
Market~14 prices \emph{lab-internal} binding authority: a precommitted process with power to delay, restrict, or cancel deployment on evidence.
That is necessary but not sufficient for a successful pause when some other market resolves NO or a live system crosses a red line.
Price institutional pause capacity separately such as via Metaculus question~44423 on whether AI safety legislation with binding effect is enacted in 2027--2028.

\paragraph{Optimistic bound.}
Let $q_{\mathrm{pause}}$ be the external governance price above (or another listed pause market traders agree on).
A deliberately optimistic upper bound on catastrophe before the tools and pause machinery are in place is schematically
\[
P(\mathrm{doom}) \;\lesssim\; 1 - \Bigl(\prod_{i \in \mathcal{P}} p_i\Bigr)\, q_{\mathrm{pause}}\, p_{15}\, p_{18},
\]
where $\mathcal{P}$ is the spine-necessary subset of Markets~1--13 and~16 along the path you treat as load-bearing for a deployment safety case (not an independent product over all eighteen rows).
This is an \emph{optimistic} bound because it assumes:
\begin{itemize}
\item a YES on market~$i$ means the \emph{tool exists}, not that it certifies the frontier system you care about;
\item benchmark-to-instance adapters (per-system scoped certificates) are available when needed;
\item shared-instrument correlation does not inflate confidence;
\item $q_{\mathrm{pause}}$ captures both Market~14-style lab authority and external legislative or coordination capacity to act on bad news.
\end{itemize}
A NO on any load-bearing row is evidence the path is not yet operationalized; it is \emph{not} by itself proof the bridge is false.


\end{authbar}
\authbarneedspace
\section{Not in the 2027 catalog}
\label{sec:appp-out-of-catalog}
\begin{authbar}{AI+GZ}
The catalog omits, on purpose:

\begin{itemize}
\item Unrestricted ``the system is safe'' with no setting.
\item A market for what residual risk people will accept (the vote is an input to Market~18).
\item Target realization or construction as a build-the-certified-class market; do not fold that into Market~14.
\item Named restorer of a control cut after damage; legal successor path (who could refuse); whether the whole specify--deploy--refuse cycle stays correctable.
Those are later construction questions, not these bars.
\item Logical induction as a forecasting instrument (Chapter~\ref{ch:towards-alignment}).
These contracts are ordinary prediction markets.
\item The retired extrapolative-correction route (Appendix~\ref{appbridge-crosswalk}).
\end{itemize}

Chapter-end ``What Would Change This View'' boxes stay qualitative and undated.
They are not these contracts, and these contracts do not replace them.

\end{authbar}
\authbarneedspace
\section{Relation to other appendices}
\label{sec:appp-relation}
\begin{authbar}{AI}
Appendix~\ref{apph-research-program} is the undated measurement program: what would have to be learned for the conditionals to bind.
This appendix is the calendar-dated public-contract layer on the same cuts.
A YES here is not a row turning green in that program, and a NO is not a disconfirmation of the corresponding chapter.

Appendix~\ref{appi-lean-proof-spine} records what follows \emph{if} named hypotheses hold.
A resolved YES is not a Lean constructor, and Lean does not import these percentage bars.

Appendix~\ref{appbridge-crosswalk} maps the same cuts to field crux names.
Appendix~\ref{appn-experimental-evidence} indexes in-repo findings that may inform prices; they cannot be the sole YES evidence.

Listing on a host platform is a separate decision.
The criteria in this appendix are the artifact either way.

\end{authbar}
